\documentclass{situationist-preprint}

\title{From Hedonism to Pluralism: A Derivation}
\msid{ETH-01-c}
\msfamily{ETH-01}
\msdiscipline{Normative and Applied Ethics}
\mstype{Alternative proof}
\mskeywords{the good; hedonism; value pluralism; dominance; reversal verdicts; method of isolation; Sidgwick; Mill; Moore}

\begin{document}
\maketitle

\begin{abstract}
ETH-01 concludes that the good is what idealized reflective endorsement converges on and that it has seven irreducible components: well-being, knowledge, achievement, relationships, fairness, beauty and autonomy. Volume~I of \emph{On the Good} (ETH-01-a) derives the plurality from seven isolation verdicts, each about two outcomes that agree exactly in six components and differ in the seventh (its Assumption~5.1). We give a proof that does not assume those verdicts. It starts where the history of the subject starts, from hedonism, and uses three premises a hedonist accepts or has conceded: a dominance principle, the recombinability of profiles, and seven reversal verdicts, each about an outcome that is better in one component, no better in any other, and better overall. Reversal verdicts require only ordinal comparisons within components, not the equal levels that isolation pairs need, and they are the verdicts on which Mill, Moore and their successors moved away from hedonism. We prove that the isolation verdicts follow, that value supervenes on no proper subset of the seven, and that any sequence of minimal repairs starting from hedonism ends at the same seven components whatever the order. A variant without recombination refutes every monotone reductive theory.
\end{abstract}

\section{Introduction}\label{s:intro}

\subsection{The principal's route}

ETH-01 holds that the good is what idealized reflective endorsement converges on, and that it is plural, with seven irreducible components. Volume~I of \emph{On the Good} (ETH-01-a) defines the good as the object of convergent idealized endorsement (its Definition~3.2), represents an outcome by a profile that assigns a level to each of the seven components, and proves that the value of outcomes does not supervene on any set of components that omits one of them (Theorem~5.2). The single substantive premise is Assumption~5.1: for each component there are two profiles that agree on all the other components and differ in value. These are the isolation verdicts, the contrastive refinement of Moore's method of isolation \citep[\S112]{Moore1903} that Definition~4.1 of ETH-01-a states.

The isolation verdicts carry the whole weight of the principal's argument, as ETH-01-a itself observes, and they have been doubted on three grounds. The first is realizability. The components are causally and sometimes constitutively connected, and outcomes that agree exactly in six components while differing in the seventh can be built, if at all, only by compensation: the deceived friend in the isolation pair for relationships must be given an item of knowledge elsewhere, of exactly the significance of the one she lacks, so that the two outcomes agree in level of knowledge (ETH-01-b, Section~3). Compensation presupposes that levels can be compared in their differences, which is more than ordinal measurement provides. The second ground is reliability. The verdicts are about exactly matched outcomes, and when the matching is exact the difference in value is often small; a small difference is easily swamped by noise in judgment, and is open to the suspicion that it reflects some difference the matching failed to remove. The third ground is dialectical. A hedonist confronted with an exactly matched pair, in which two outcomes agree in everyone's well-being and differ only in, say, knowledge, will say that the pair is one of which she has no confident verdict; and the verdicts that moved the tradition away from hedonism were not of this form.

\subsection{Thesis}

We give a proof of the plurality of the good that does not assume Assumption~5.1. Its premises are these. Definition~3.2 of ETH-01-a, the account of the good as the object of convergent idealized endorsement, is retained, together with the representation of outcomes by profiles. In place of the isolation verdicts we use a dominance principle: if one outcome is at least as high as another in every component, it is at least as good. We use a recombination assumption: the levels of different components can be combined freely, which is already presupposed by the formal model of ETH-01-a, where profiles are arbitrary assignments of levels. And we use seven \emph{reversal verdicts}: for each component, there is an outcome that is higher in that component, no higher in any other, and better overall than its rival. A reversal verdict says that a gain in one component outweighs some loss, possibly very small, in others.

From these premises we prove that the isolation verdicts hold (Lemma~\ref{l:iso}), so that the principal's theorem follows, and we give a direct proof of it (Theorem~\ref{t:plural}). We then show that if recombination is given up, the reversal verdicts alone still refute every reductive theory that makes value an increasing function of a proper subset of the components (Proposition~\ref{p:norecomb}), which includes every monistic theory that has been seriously proposed. Finally we show that the derivation from hedonism has a determinate end: any sequence of minimal repairs that adds a component whenever a reversal verdict refutes the current theory reaches the full list of seven, in exactly as many steps as there are components missing at the start, whatever the order in which the reversal verdicts are met (Proposition~\ref{p:repair}).

The premise avoided is Assumption~5.1 of ETH-01-a, taken as a premise. The proof derives it, so it is not denied; but the verdicts the proof needs are of a different form. They require only that the outcome favored in the tested component be no better in the others, which is an ordinal condition, easily secured by stipulating losses rather than calibrated compensations; and they are the verdicts on which the history of value theory, from Mill onward, has in fact turned.

\subsection{Plan}

Section~\ref{s:history} tells the history of the passage from hedonism to pluralism as a sequence of reversal verdicts. Section~\ref{s:frame} sets out the framework. Section~\ref{s:deriv} gives the derivation. Section~\ref{s:cases} states the seven reversal verdicts, and Section~\ref{s:starts} applies the derivation to starting points other than hedonism. Section~\ref{s:premise} explains what the avoided premise did and what its avoidance costs. Section~\ref{s:obj} answers objections, and Section~\ref{s:not} states what the result does not establish.

\section{From hedonism to pluralism: the history}\label{s:history}

The passage from hedonism to pluralism in modern value theory can be told as a sequence of cases, each of which presented an outcome or a life that was worse in pleasure and better overall, and each of which led to a revision of the theory.

\paragraph{Sidgwick.} The most careful statement of hedonism is Sidgwick's. He characterized a person's good on the whole as what she would desire if all the consequences of the alternatives open to her were accurately foreseen and adequately realized in imagination \citep[I.ix.3]{Sidgwick1907}, an idealized-desire account that is the ancestor of Definition~3.1 of ETH-01-a. He then argued, in the chapter on ultimate good \citep[III.xiv]{Sidgwick1907}, that when we reflect on the things other than pleasure that are commonly held good, such as knowledge, virtue and beauty, we find that we value them only so far as they conduce to desirable consciousness. Sidgwick's argument is therefore an argument from the absence of reversal verdicts: it claims that idealized reflection never prefers an outcome lower in pleasure on account of something else. The claim is testable, and the history that follows is a series of tests.

\paragraph{Mill.} Mill had already given the first counterexample, in the form of a reversal. It is better, he held, to be a human being dissatisfied than a pig satisfied, and better to be Socrates dissatisfied than a fool satisfied \citep[ch.~II]{Mill1861}. The comparison is between two lives of which the first contains less contentment and more of the exercise of the higher faculties, and the verdict is that the first is better. Mill presented the verdict as a refinement of hedonism, distinguishing pleasures by quality as well as quantity. \citet[\S\S46--48]{Moore1903} replied that if pleasures differ in value by something other than their quantity as pleasures, then something other than pleasure is good, and hedonism has been abandoned. Mill's verdict, in the terms of the present manuscript, is a reversal verdict for knowledge and the exercise of capacities; his revision admitted a second dimension of value while retaining the name of the first.

\paragraph{Moore and Rashdall.} Moore's own list of intrinsic goods, of which the greatest are personal affection and the appreciation of beauty \citep[\S113]{Moore1903}, was reached by the method of isolation, but the cases he used to persuade were often reversals. His comparison of a beautiful and an ugly world, neither ever seen \citep[\S50]{Moore1903}, was put against Sidgwick's claim that beauty apart from contemplation has no value. \citet{Rashdall1907} defended an ideal utilitarianism in which virtue, knowledge and culture, as well as pleasure, are elements of the good, again on the ground that reflection prefers outcomes containing them even at some cost in pleasure. \citet{Ross1930} added the apportionment of pleasure to virtue, a distributive good of the outcome as a whole, defended by a comparison of two worlds alike in their virtue and their pleasure, in one of which the virtuous are happy and in the other the vicious; his case is an isolation case rather than a reversal, but the distributive literature that followed restated such cases as reversals.

\paragraph{The twentieth century.} \citet[pp.~42--45]{Nozick1974} gave the experience machine, which ETH-05-a restates as a comparison between experiential twins and uses to show that well-being includes contact with reality. The machine is a reversal in its original form, since the life in the machine may be stipulated to be more pleasant than the life outside it. \citet{Griffin1986} and \citet{Parfit1984} drew from such cases the objective-list theory of well-being, which places knowledge, achievement and deep personal relationships among the things that make a life go well. The prioritarian and egalitarian literature that followed \citet{Diamond1967} and \citet{Parfit1997} made distributive features of outcomes a central topic, and in every version the cases were of the same form: an outcome lower in total well-being, and better distributed, judged better.

\paragraph{Recent hedonism.} Hedonism has been revived in the last two decades, in forms designed to absorb the reversal verdicts. \citet{Feldman2004} defends attitudinal hedonism, on which pleasure is an attitude toward a state of affairs, and considers versions on which the value of such a pleasure is adjusted for the truth of its object and for the worth of the object, so that pleasure taken in what is not the case, or in what is base, counts for less. \citet{Crisp2006} defends a hedonism about well-being on which enjoyment is a determinable whose determinates differ in value, and \citet{Bramble2016} argues that the experience machine is answered by noting that the pleasures of real achievement and relationship are different pleasures. Each of these positions concedes, in the terms of Moore's reply to Mill, that the value of a pleasure depends on something other than its pleasantness: on the truth of its object, on its worth, or on its source. They differ from pluralism in how the dependence is represented, not in whether it exists, and the argument below shows that the representation does not matter.

The history therefore has a structure. A theory makes value a function of some features. A reversal verdict presents an outcome that is no better in those features and better overall. The theory is revised by adding a feature. The present manuscript turns this structure into a proof.

\section{Framework}\label{s:frame}

We use the representation of ETH-01-a. An outcome is represented by a profile $x$ that assigns to each of the seven components $i \in \mathcal{C}$ a level $x_i$, and $V$ is the value assigned to profiles by idealized reflective endorsement in the sense of Definition~3.2 of ETH-01-a. Levels within a component are assumed to be ordered: of two outcomes, one has at least as much knowledge as the other, or more fairness, and so on. No assumption is made about the comparability of differences in level, or about comparisons between levels of different components.

\begin{definition}\label{d:dom}
A profile $x$ \emph{weakly dominates} a profile $y$ if $x_j \geq y_j$ for every component $j$. For a component $i$, $x$ is \emph{no better than $y$ outside $i$} if $x_j \leq y_j$ for every component $j \neq i$.
\end{definition}

\begin{definition}\label{d:reversal}
A \emph{reversal pair} for a component $i$ is a pair of profiles $(x, y)$ such that $x_i > y_i$, $x$ is no better than $y$ outside $i$, and $V(x) > V(y)$.
\end{definition}

A reversal pair presents an outcome that is better in $i$, no better in anything else, and better. Its rival may be better in every other component, and the size of its advantage there is not specified.

\begin{principle}[Dominance]\label{p:dom}
If $x$ weakly dominates $y$, then $V(x) \geq V(y)$.
\end{principle}

\begin{assumption}[Recombination]\label{a:recomb}
For any profiles $x$ and $y$ and any set $S$ of components, the assignment that agrees with $x$ on $S$ and with $y$ on the other components is a profile, that is, the profile of some possible outcome.
\end{assumption}

\begin{assumption}[Reversal verdicts]\label{a:rev}
For each of the seven components there is a reversal pair.
\end{assumption}

Principle~\ref{p:dom} says that more of a component of the good, with no less of any other, is never worse. It is the weakest principle that connects the components to value at all, and it is accepted by every monist for the components she recognizes: for the hedonist, more well-being with no less of anything else is not worse, and she is indifferent to the rest. Assumption~\ref{a:recomb} is part of the formal model of ETH-01-a, in which profiles are arbitrary functions from components to levels; we make it explicit because Section~\ref{s:deriv} shows what can be proved without it. Assumption~\ref{a:rev} is the substantive premise, and Section~\ref{s:cases} states the seven verdicts.

We also need the notion of a reductive theory.

\begin{definition}\label{d:reductive}
A \emph{reductive theory with base $S$}, for a set $S \subseteq \mathcal{C}$, holds that $V(x) = f(x|_S)$ for some function $f$ of the levels of the components in $S$. It is \emph{monotone} if $f$ is weakly increasing in each of its arguments.
\end{definition}

A reductive theory with base $S$ is true just in case $V$ supervenes on $S$. Hedonism, as a theory of the good of outcomes, is a monotone reductive theory whose base is well-being, with well-being measured hedonically. Every monism that has been proposed is a monotone reductive theory with a one-element base: no monist holds that more of her good can make an outcome worse.

\section{The derivation}\label{s:deriv}

\begin{lemma}\label{l:iso}
Given Principle~\ref{p:dom} and Assumption~\ref{a:recomb}, every reversal pair for a component $i$ yields an isolation pair for $i$: two profiles that agree on every component other than $i$ and differ in value.
\end{lemma}

\begin{proof}
Let $(x, y)$ be a reversal pair for $i$. By Assumption~\ref{a:recomb} there is a profile $z$ that agrees with $x$ on $i$ and with $y$ on every other component. For each $j \neq i$, $z_j = y_j \geq x_j$, since $x$ is no better than $y$ outside $i$; and $z_i = x_i$. So $z$ weakly dominates $x$, and by Principle~\ref{p:dom}, $V(z) \geq V(x) > V(y)$. The profiles $z$ and $y$ agree on every component other than $i$, and $V(z) \neq V(y)$.
\end{proof}

With Assumption~\ref{a:rev}, Lemma~\ref{l:iso} yields Assumption~5.1 of ETH-01-a, and Theorem~5.2 of ETH-01-a then gives plurality. We give the direct proof, which shows the role of each premise.

\begin{theorem}\label{t:plural}
Given Principle~\ref{p:dom} and Assumptions~\ref{a:recomb} and~\ref{a:rev}, $V$ does not supervene on any set of components that omits at least one of the seven.
\end{theorem}

\begin{proof}
Let $S$ omit $i$, and let $(x, y)$ be a reversal pair for $i$, which exists by Assumption~\ref{a:rev}. Let $z$ agree with $x$ on $S$ and with $y$ elsewhere, which is a profile by Assumption~\ref{a:recomb}. Since $i \notin S$, every component in $S$ is a component other than $i$, so $z_j = x_j \leq y_j$ for $j \in S$, and $z_j = y_j$ for $j \notin S$; hence $y$ weakly dominates $z$, and $V(z) \leq V(y)$ by Principle~\ref{p:dom}. If $V$ supervened on $S$, then since $z$ and $x$ agree on $S$, $V(z) = V(x) > V(y)$, a contradiction.
\end{proof}

\begin{corollary}\label{c:monism}
Given the same premises, the value of an outcome is not a function of any single component; in particular, hedonism and every theory on which the value of an outcome is a function of well-being alone are false.
\end{corollary}

\begin{proof}
A single component is a set that omits six of the seven.
\end{proof}

Corollary~\ref{c:monism} reproduces Corollaries~5.3 and~5.4 of ETH-01-a. Its application to hedonism is overdetermined: ETH-05-a refutes hedonism about well-being by showing that well-being includes contact with reality (its Theorem~3.2), and Corollary~\ref{c:monism} refutes the more general welfarist thesis that the good of outcomes is a function of well-being however measured.

Recombination is the premise most open to doubt (Objection~\ref{o:recomb} below), and it is not needed to refute the theories that have actually been proposed.

\begin{proposition}\label{p:norecomb}
Given Assumption~\ref{a:rev} alone, no monotone reductive theory whose base omits a component is true.
\end{proposition}

\begin{proof}
Let the theory have base $S$ omitting $i$, so that $V(x) = f(x|_S)$ with $f$ weakly increasing. Let $(x, y)$ be a reversal pair for $i$. Every component in $S$ is other than $i$, so $x_j \leq y_j$ for each $j \in S$, and since $f$ is weakly increasing, $f(x|_S) \leq f(y|_S)$, that is, $V(x) \leq V(y)$. This contradicts $V(x) > V(y)$.
\end{proof}

The proposition uses neither Dominance nor Recombination, and it needs only one reversal pair for each component omitted. Its conclusion is weaker than Theorem~\ref{t:plural} in one respect: it leaves open that $V$ supervenes on a proper subset of the components through a function that decreases in some of them. Such a theory would hold that more of some component of the good, everything else in its base equal, can make an outcome worse. No such theory has been proposed, and with Dominance and Recombination it is excluded by Theorem~\ref{t:plural}.

The history of Section~\ref{s:history} was a sequence of revisions. We now show that its endpoint does not depend on the order in which the cases arrived.

\begin{definition}\label{d:repair}
A \emph{repair sequence} is a sequence of bases $S_0 \subseteq S_1 \subseteq \cdots$ such that each $S_{k+1}$ is $S_k \cup \{i\}$ for some component $i \notin S_k$ for which a reversal pair exists, and which stops at $S_n$ only if no component outside $S_n$ has a reversal pair.
\end{definition}

A repair sequence models a theorist who holds a monotone reductive theory with base $S_k$, meets a reversal pair for a component outside it, which refutes the theory by Proposition~\ref{p:norecomb}, and makes the minimal revision of adding that component and nothing else.

\begin{proposition}\label{p:repair}
Given Assumption~\ref{a:rev}, every repair sequence starting from a base $S_0$ ends at $\mathcal{C}$, after exactly $|\mathcal{C}\setminus S_0|$ steps, whatever the order in which components are added.
\end{proposition}

\begin{proof}
Each step adds exactly one component not already in the base, so the sequence has at most $|\mathcal{C}\setminus S_0|$ steps. By Assumption~\ref{a:rev}, every component has a reversal pair, so the sequence cannot stop while some component lies outside the base; it stops only at $\mathcal{C}$. Hence it reaches $\mathcal{C}$ after adding each of the $|\mathcal{C}\setminus S_0|$ missing components once.
\end{proof}

Starting from hedonism, $S_0$ is the base consisting of well-being, and the sequence has six steps. A preliminary step occurs within the first component, when hedonism about well-being gives way, through the experience machine, to a conception of well-being that includes contact with reality (ETH-05-a); that step changes how the first component is measured, not which components there are. The order in which the history took the remaining steps, roughly knowledge, beauty and relationships with Mill and Moore, fairness with the distributive literature that followed Ross, achievement and autonomy with the objective-list theorists and the liberal tradition, is one of the $6! = 720$ orders, and Proposition~\ref{p:repair} shows that it made no difference to the destination.

Finally, we record the exact logical relation between the two kinds of verdict.

\begin{proposition}\label{p:equiv}
Given Principle~\ref{p:dom}, every isolation pair for a component yields a reversal pair for it. Given Principle~\ref{p:dom} and Assumption~\ref{a:recomb}, a component has a reversal pair if and only if it has an isolation pair.
\end{proposition}

\begin{proof}
Let $z$ and $y$ agree on every component other than $i$ and differ in value. They differ in $i$, since otherwise they would have the same profile; label them so that $z_i > y_i$. Then $z$ weakly dominates $y$, so $V(z) \geq V(y)$ by Principle~\ref{p:dom}, and since their values differ, $V(z) > V(y)$. As $z$ agrees with $y$ outside $i$, it is no better than $y$ outside $i$, and $(z, y)$ is a reversal pair. The converse, given Assumption~\ref{a:recomb}, is Lemma~\ref{l:iso}.
\end{proof}

Proposition~\ref{p:equiv} shows that, with Dominance and Recombination in place, the existential content of Assumption~\ref{a:rev} and that of Assumption~5.1 of ETH-01-a coincide: an isolation pair is a reversal pair with no loss. The two proofs differ in which pairs they offer to idealized judgment as witnesses, and Section~\ref{s:premise} explains why the difference matters. Without Recombination they differ in content as well: Proposition~\ref{p:norecomb} uses reversal pairs that need not be convertible into isolation pairs, and so it applies where outcomes exactly matched in six components do not exist.

\section{The seven reversal verdicts}\label{s:cases}

For each component we state a reversal pair. In each, the first outcome $x$ is higher in the tested component, the losses it suffers in other components are stipulated, and the verdict is that $x$ is better than $y$. The losses can be made small, but they need not be calibrated against anything: it suffices that $x$ is no better than $y$ in any component other than the one tested.

\paragraph{Well-being.} A community can relieve a chronic and severe pain suffered by many of its members by building a clinic on a hillside whose rock formations, seen by no one, are of some beauty; the formations are destroyed. The relief is the only effect on well-being; nothing else is gained. The outcome with the clinic is better. No one disputes the verdict, and it is included only for completeness: well-being is the one component that every party to the debate recognizes.

\paragraph{Knowledge.} A society funds a program of basic research that will establish, among other things, the age and composition of the early universe. The knowledge will have no practical application, and the funds come from a modest reduction in public amenities, so that the members' lives are slightly less comfortable; nothing else is gained. The outcome with the research is better. This is the form of Mill's verdict about Socrates, stated for an outcome rather than a life, and with the gain confined to knowledge.

\paragraph{Achievement.} A composer spends a decade on a difficult work and completes it, at the cost of some ease and comfort she would otherwise have had; the work is never performed and gives no one pleasure, and she learns nothing she would not otherwise have learned. The outcome in which she completes it is better than the one in which she lives the easier life. \citet{Bradford2015} defends the value of achievement, understood as the competent causation of a difficult outcome, on cases of this form.

\paragraph{Relationships.} Two people who would otherwise remain acquaintances become close friends. The friendship costs each some ease, since each comes to share the other's troubles, and its end, with the death of one, brings the other grief; in total, their lives contain somewhat less pleasure than they would have without it, and nothing else is gained. The outcome with the friendship is better. The verdict is the old one that it is better to have loved and lost, and it is the verdict on which Moore placed personal affection among the greatest goods \citep[\S113]{Moore1903}.

\paragraph{Fairness.} A single indivisible good is to be allocated between two people with equal claims. In one outcome it is allocated by a fair lottery; in the other, it is given to the one who would benefit slightly more, because her name comes first in the alphabet. The total well-being in the first outcome is slightly lower in expectation, and no other component differs. The first outcome is better. This is the verdict of \citet{Diamond1967}, and \citet{Broome1990} explains it: when claims are equal and the good indivisible, a lottery is the only way of giving each claim its due.

\paragraph{Beauty.} The hillside formation of the first pair is preserved instead, at a very small cost: a road is built around it, adding a minute to the journeys of those who use it, and no one ever sees the formation. The outcome with the formation is better. This verdict is the weakest of the seven, as the isolation verdict for beauty is in ETH-01-a, and \citet[III.xiv]{Sidgwick1907} denied it. But the verdict requires only that the value of unseen beauty exceed some loss, and the loss can be made as small as we please. Sidgwick's argument concerned what it is rational to aim at, and Moore's reply, that unseen beauty has some value though little \citep[\S\S50, 113]{Moore1903}, is all the verdict needs.

\paragraph{Autonomy.} A young person chooses her occupation for herself, and chooses a little less well, for her own well-being, than a benevolent adviser would have chosen for her by arranging her circumstances without her knowledge; nothing else differs. The outcome in which she chooses is better. This is Mill's verdict in the chapter of \emph{On Liberty} on individuality \citep[ch.~III]{Mill1859}, and \citet{Raz1986} defends it as the value of personal autonomy, given an adequate range of options.

\begin{remark}\label{r:ordinal}
Each pair is constructed by stipulating that the first outcome is worse in one or more components other than the tested one, and otherwise the same. The construction needs only the ordering of levels within each component: a slightly less comfortable life is lower in well-being, a destroyed formation lowers the level of beauty. It needs no judgment that a loss in one component is exactly matched by a compensating gain in the same component elsewhere.
\end{remark}

\section{Other starting points}\label{s:starts}

The derivation has been told from hedonism, because that is where the history starts. Proposition~\ref{p:repair} applies to every starting base, and three other starting points have been historically important.

\paragraph{Desire satisfaction.} Desire-satisfaction theories of well-being were introduced in part to absorb the experience machine: a person who wants to write novels does not get what she wants by having the experience of writing them. As a theory of the good of outcomes, the view makes value a function of the satisfaction of desires, and its base, in the terms of the profile, is well-being measured by desire satisfaction. \citet{Heathwood2006} has argued that, on its most defensible version, desire satisfactionism converges with a form of hedonism, since the desires that matter are intrinsic desires for experiences. On either version, the base is a single component, and the repair sequence is the one from hedonism. Reversal verdicts for fairness, beauty and autonomy are especially telling against it, since in each the favored outcome may satisfy fewer desires: the lottery frustrates the desire of the candidate who would have received the good, the unseen formation is the object of no one's desire, and the adviser's arrangement satisfies more of the young person's desires than her own choice does.

\paragraph{Perfectionism.} \citet{Hurka1993} defends a perfectionism on which the good is the development of human nature, centrally of theoretical and practical rationality, which yields the value of knowledge and achievement. Its base, in the terms of the profile, is knowledge and achievement, and the repair sequence from it has five steps. The first, which perfectionists themselves have often taken, is the addition of well-being, on the reversal verdict that the relief of severe suffering at a small cost in knowledge is better. Relationships, fairness, beauty and autonomy follow on the verdicts of Section~\ref{s:cases}. Perfectionists have sometimes argued that autonomy and relationships are themselves forms of the exercise of practical rationality; if so, the base is larger at the start, and the sequence shorter, but its endpoint is the same.

\paragraph{Ross.} Ross's list of intrinsic goods, virtue, pleasure, knowledge and the apportionment of pleasure to virtue \citep{Ross1930}, starts from a base that is both narrower and broader than the seven: narrower, since it lacks achievement, relationships, beauty and autonomy, and broader, since it contains virtue. Repairs add components; they do not remove them. The removal of virtue requires a different argument: that virtue has no reversal verdict relative to the seven, so that two outcomes alike in all seven do not differ in value on account of virtue. That is the argument ETH-01-d gives, with its Proposition~1.1, which shows that value supervenes on a set of components exactly when no feature passes the isolation test relative to it. The derivation from hedonism and the closure argument are therefore complementary: the first shows that the list must contain at least the seven, and the second, which belongs to Volume~VI, that it need contain no more.

\section{What the avoided premise did}\label{s:premise}

Assumption~5.1 of ETH-01-a asserts, for each component, the existence of two outcomes that agree exactly on the other six and differ in value. The present proof derives these pairs (Lemma~\ref{l:iso}) from reversal pairs, Dominance and Recombination. It is worth being precise about what the substitution gains and what it costs.

\paragraph{What is gained.} First, the verdicts required are ordinal in construction (Remark~\ref{r:ordinal}). An isolation pair in which the concrete outcomes differ in more than one respect must be completed by compensation, and compensation within a component presupposes that equal differences in level can be identified (ETH-01-b, Remark~3.1). A reversal pair needs no compensation: if the tested outcome suffers an incidental loss in some other component, the pair remains a reversal pair. Second, the verdicts are robust to the error that most often affects isolation judgments, the importation of an unstated difference. If a judge, contemplating the research program, imagines that its knowledge will eventually be useful, she imports a difference in well-being; but in a reversal pair the stipulated loss in well-being can be enlarged until any such imported gain is outweighed, and the verdict persists for losses that are small. Third, the verdicts are those on which the tradition actually moved. Mill's Socrates, Moore's beautiful world as put against Sidgwick, Nozick's machine and Diamond's lottery are reversal cases. A hedonist who rejects the isolation verdicts as too faint to assess is confronted, in the present proof, with verdicts her own tradition conceded, together with a principle, Dominance, that she accepts.

\paragraph{What it costs.} The proof needs Dominance and Recombination, which the principal's proof does not need as separate premises. Recombination is presupposed by the formal model of ETH-01-a, so the cost is one of explicitness; and Proposition~\ref{p:norecomb} shows that the refutation of every monotone reductive theory survives without it. Dominance is a genuine additional premise. It is weak, and it is accepted by every monist for her own good, but it would be rejected by a theory on which more of some component can, with nothing else lower, make an outcome worse. We consider this possibility in Objection~\ref{o:dom}. The particular reversal verdicts of Section~\ref{s:cases} are also stronger than particular isolation verdicts: each asserts that a gain in one component outweighs a specified loss in another, and not merely that it makes some difference. As existential claims, by Proposition~\ref{p:equiv}, the two kinds of verdict coincide given Dominance and Recombination; as evidence, each reversal verdict commits the judge to more. Their advantage is that the loss may be stipulated as small as one likes, so that a reversal verdict with a negligible loss asserts little more than an isolation verdict, while retaining the ordinal construction. The premise avoided is therefore Assumption~5.1 taken as a report of verdicts on exactly matched outcomes; what replaces it is a set of verdicts on outcomes matched only in the ordinal sense, from which the exactly matched pairs are constructed rather than judged.

\paragraph{What is shared.} Both proofs rely on Definition~3.2 of ETH-01-a, the account of the good as the object of convergent idealized endorsement, and both use the profile representation. Neither establishes completeness. The two proofs therefore support the conclusion of ETH-01 from different evidence about pairs of outcomes, through the same account of what makes an outcome better.

\section{Objections}\label{s:obj}

\begin{objection}[Dominance]\label{o:dom}
Dominance fails for several components. More knowledge of trivial facts does not make an outcome better, and may make it worse by crowding out more valuable pursuits. More options can make a life worse, as \citet{Dworkin1982} argued, so more autonomy is not always better. Some relationships are destructive.
\end{objection}

\begin{reply}
Each example confuses a level of a component with a quantity of something that bears on it. The knowledge component records the value-relevant level of knowledge, which weights items by their significance \citep{Hurka1993}; trivial knowledge adds little or nothing, and does not lower the level. Whatever crowding-out it causes lowers some other component, so it is not a case of more of one with no less of the others. Dworkin's argument concerns the number of options, not the degree to which a life is self-directed, and its examples are cases in which more options reduce self-direction, by imposing burdens of choice or exposing the agent to pressure; they lower the level of autonomy. A destructive relationship is low in the relationship component, as that component is understood in ETH-01: what is measured is the presence of genuine mutual care, not the number or intensity of entanglements. Dominance is therefore not a substantive claim about the effects of having more of things; it is the claim that the components are components of the good.
\end{reply}

\begin{objection}[Recombination]\label{o:recomb}
Some combinations of levels are impossible. There can be no high level of achievement with no knowledge, and no relationships without some well-being. Recombination is false, and Theorem~\ref{t:plural} rests on it.
\end{objection}

\begin{reply}
Recombination requires only that the levels used in the proof can be combined, and the proof combines levels that are close to those of the reversal pair itself: the profile $z$ in the proof of Theorem~\ref{t:plural} takes the levels of $x$ on some components and those of $y$ on the rest, and since the two outcomes differ by small stipulated amounts, $z$ is a minor variant of each. But the reply does not depend on this. Proposition~\ref{p:norecomb} refutes every monotone reductive theory with a base that omits a component without using Recombination at all. The only theories it leaves standing are those that make value a non-monotone function of a proper subset of the components, and those are excluded by Dominance once Recombination holds even locally.
\end{reply}

\begin{objection}[Qualitative hedonism]
The reversal verdicts are compatible with hedonism if pleasures differ in value by quality, as \citet{Mill1861} held, or if the value of a pleasure is adjusted for the truth or worth of its object, as in the theories considered by \citet{Feldman2004}. The research program produces pleasures of a higher quality; the fair lottery produces pleasure in fairness; and so on.
\end{objection}

\begin{reply}
Each revision makes the value of an outcome depend on features other than the quantity of pleasure it contains: on whether a pleasure is taken in knowledge, whether its object is real, whether it is taken in something worthy. The reversal pairs are stated with these features in view. In the pair for knowledge, the research may produce no pleasure in anyone, since its results are recorded and not contemplated; in the pair for beauty, no one sees the formation. A truth-adjusted or desert-adjusted hedonism cannot assign the first outcome higher value in such pairs without letting value depend on the knowledge or the beauty itself. That is Moore's point against Mill \citep[\S48]{Moore1903}, and it applies to every hedonism that absorbs the reversals: the result is a theory whose base includes the components, under the name of qualities of pleasure. Proposition~\ref{p:repair} shows that the destination of the repairs does not depend on what the added components are called.
\end{reply}

\begin{objection}[Why seven]
The derivation stops at seven because the list it starts from has seven entries. A repair sequence over a longer list of candidates would end with a longer list.
\end{objection}

\begin{reply}
The objection is correct about the derivation and does not touch the conclusion it supports. Proposition~\ref{p:repair} shows that repairs from hedonism reach every component for which there is a reversal verdict, and Assumption~\ref{a:rev} asserts such verdicts for the seven. That no further candidate has a reversal verdict is the completeness thesis, which identifies the closure of the list with the supervenience of value on all seven components together. It is argued in Volume~VI of ETH-01 and examined in ETH-01-d, which tests the principal further candidates (virtue, life, religion, play, equality and rights) and finds that none passes the isolation test relative to the seven. A candidate with a reversal verdict would, by Lemma~\ref{l:iso} applied to the enlarged list, pass that test, so the candidates ETH-01-d rejects have no reversal verdicts either.
\end{reply}

\section{What the result does not establish}\label{s:not}

First, the result does not establish that the list is complete. Theorem~\ref{t:plural}, like Theorem~5.2 of ETH-01-a, shows that the seven components are necessary; that value supervenes on them together is the thesis of Volume~VI.

Second, it does not establish the account of the good on which it relies. Definition~3.2 of ETH-01-a, on which the good is what idealized reflective endorsement converges on, is assumed; Sidgwick's idealized-desire criterion is its ancestor, and MTE-01 argues for the general claim that moral facts are natural facts about idealized reactions. The present proof avoids Assumption~5.1 of ETH-01-a, not Definition~3.2.

Third, it does not establish how the components are weighed. A reversal verdict asserts that a particular gain outweighs a particular loss, and the seven verdicts of Section~\ref{s:cases} are each compatible with a very wide range of trade-offs between components. Nothing follows about whether the components are commensurable in general.

Fourth, Proposition~\ref{p:norecomb} without Recombination refutes only monotone reductive theories. A theory on which more of some component can, other things in its base equal, make an outcome worse is refuted only with Dominance and Recombination.

\section{Conclusion}

The plurality of the good can be derived without the isolation verdicts of ETH-01-a. From a dominance principle that every monist accepts for her own good, the recombinability of profiles that the formal model of ETH-01-a presupposes, and seven reversal verdicts of the kind on which Mill, Moore, Nozick and Diamond moved the subject away from hedonism, the isolation verdicts follow, and with them the conclusion that value supervenes on no set of components that omits one of the seven. Without recombination, the reversal verdicts alone refute every monotone reductive theory. The passage from hedonism to pluralism is a sequence of minimal repairs, and every such sequence ends at the same seven components.

\begin{thebibliography}{99}\small
\bibitem[Bradford(2015)]{Bradford2015} Bradford, G. (2015). \emph{Achievement}. Oxford University Press.
\bibitem[Bramble(2016)]{Bramble2016} Bramble, B. (2016). A new defense of hedonism about well-being. \emph{Ergo}, 3(4).
\bibitem[Broome(1990)]{Broome1990} Broome, J. (1990). Fairness. \emph{Proceedings of the Aristotelian Society}, 91, 87--101.
\bibitem[Crisp(2006)]{Crisp2006} Crisp, R. (2006). \emph{Reasons and the Good}. Clarendon Press.
\bibitem[Diamond(1967)]{Diamond1967} Diamond, P.~A. (1967). Cardinal welfare, individualistic ethics, and interpersonal comparison of utility: Comment. \emph{Journal of Political Economy}, 75(5), 765--766.
\bibitem[Dworkin(1982)]{Dworkin1982} Dworkin, G. (1982). Is more choice better than less? \emph{Midwest Studies in Philosophy}, 7, 47--61.
\bibitem[Feldman(2004)]{Feldman2004} Feldman, F. (2004). \emph{Pleasure and the Good Life: Concerning the Nature, Varieties, and Plausibility of Hedonism}. Clarendon Press.
\bibitem[Griffin(1986)]{Griffin1986} Griffin, J. (1986). \emph{Well-Being: Its Meaning, Measurement and Moral Importance}. Clarendon Press.
\bibitem[Heathwood(2006)]{Heathwood2006} Heathwood, C. (2006). Desire satisfactionism and hedonism. \emph{Philosophical Studies}, 128(3), 539--563.
\bibitem[Hurka(1993)]{Hurka1993} Hurka, T. (1993). \emph{Perfectionism}. Oxford University Press.
\bibitem[Mill(1859)]{Mill1859} Mill, J.~S. (1859). \emph{On Liberty}. John W. Parker and Son.
\bibitem[Mill(1861)]{Mill1861} Mill, J.~S. (1861). Utilitarianism. \emph{Fraser's Magazine}, 64; in book form, Parker, Son and Bourn, 1863.
\bibitem[Moore(1903)]{Moore1903} Moore, G.~E. (1903). \emph{Principia Ethica}. Cambridge University Press.
\bibitem[Nozick(1974)]{Nozick1974} Nozick, R. (1974). \emph{Anarchy, State, and Utopia}. Basic Books.
\bibitem[Parfit(1984)]{Parfit1984} Parfit, D. (1984). \emph{Reasons and Persons}. Clarendon Press.
\bibitem[Parfit(1997)]{Parfit1997} Parfit, D. (1997). Equality and priority. \emph{Ratio}, 10(3), 202--221.
\bibitem[Rashdall(1907)]{Rashdall1907} Rashdall, H. (1907). \emph{The Theory of Good and Evil}. Clarendon Press.
\bibitem[Raz(1986)]{Raz1986} Raz, J. (1986). \emph{The Morality of Freedom}. Clarendon Press.
\bibitem[Ross(1930)]{Ross1930} Ross, W.~D. (1930). \emph{The Right and the Good}. Clarendon Press.
\bibitem[Sidgwick(1907)]{Sidgwick1907} Sidgwick, H. (1907). \emph{The Methods of Ethics} (7th ed.). Macmillan.
\end{thebibliography}

\end{document}
